% C6 — A mechanized, non-circular renormalization-group interface for
% Wilson lattice gauge correlators.  Charter: docs/C6-BRIDGE-CHARTER.md
% (Amendments 1-9; Amendment 9 = the owner closure order fixing title
% direction, section structure, and the mandatory prohibitions).
% REGIME: tricotomy; countermodel history presented as adversarial
% design; NO scores, NO desk language, NO commit chronology; freeze at
% d75d8952 (this manuscript added in a strictly editorial later commit).
\documentclass[11pt]{article}
\usepackage[margin=1.1in]{geometry}
\usepackage{amsmath,amssymb,amsthm}
\usepackage{booktabs}
\usepackage[colorlinks=true,linkcolor=blue,urlcolor=blue,citecolor=blue]{hyperref}

\newtheorem{theorem}{Theorem}[section]
\newtheorem{lemma}[theorem]{Lemma}
\newtheorem{proposition}[theorem]{Proposition}
\newtheorem{remark}[theorem]{Remark}
\theoremstyle{definition}
\newtheorem{definition}[theorem]{Definition}
\newtheorem{interface}[theorem]{Interface}
\newtheorem{countermodel}[theorem]{Countermodel}

% long monospace artifact names offer no hyphenation points;
% let paragraphs stretch rather than overflow the margin
\emergencystretch=3em
\tolerance=2000

\newcommand{\lean}[1]{\texttt{#1}}
\newcommand{\Torus}[1]{\Lambda_{#1}}
\newcommand{\Conf}[1]{\mathcal{A}_{#1}}
\newcommand{\SU}{\mathrm{SU}}
\newcommand{\dist}{\mathrm{dist}}
\newcommand{\holo}{\mathrm{hol}}
\newcommand{\Cov}{\mathcal{C}}
\newcommand{\muW}{\mu^{\mathrm{W}}}
\newcommand{\kT}{k_{\mathrm{T}}}

\title{A Mechanized, Non-Circular Renormalization-Group Interface\\
for Wilson Lattice Gauge Correlators}
\author{Lluis Eriksson\thanks{This manuscript and parts of the formal
development were produced with AI assistance under the repository's
audit protocol.  Every statement labelled \emph{exact} below is
machine-checked in Lean~4 against a pinned Mathlib with a clean axiom
oracle; countermodel analyses that are not formalized are labelled
paper-level where they occur.  Repository:
\url{https://github.com/lluiseriksson/THE-ERIKSSON-PROGRAMME}.}}
\date{July 2026}

\begin{document}
\maketitle

\begin{abstract}
For $\SU(N_c)$ Wilson lattice gauge theory on $d$-dimensional periodic
tori ($d \ge 2$), we present a machine-checked (Lean~4, pinned Mathlib)
renormalization-group \emph{interface} for two-plaquette truncated
correlators, in which every structural ingredient is a theorem rather
than a postulate: the scale transformation is a concrete decimation
map, defined once --- measurable, local, gauge-covariant, and
probability-preserving; the effective measures are its literal iterated
pushforwards of the Wilson Gibbs measure; the multiscale decomposition
of the correlator is proved by telescoping, never carried as data; the
terminal scale of the decomposition is a fixed index $\kT(n) = n$ with
typed range $1 \le \kT(n) \le n$, which excludes, in the type, the
circular depth-zero layer in which an infrared clause would hypothesize
the bound being sought; the conditional decay conclusion is stated in
the physical distance $2^n u$ with a single constant pair $(C, m)$
quantified before every torus base and depth; and the terminal
observable is \emph{operationally support-certified}: the infrared
object consumed by the interface equals a base-measure integral of an
explicitly composed pullback observable whose dependence is contained
in a transported support set, for which the separation lower bound
$2^n(2u) - (2^n + 1)$, strictly positive on the whole interface window,
is proved.  The design is deliberately adversarial: four natural naive
formulations are presented together with the explicit countermodels
that defeat them --- scalar relabeling of known decay, sink flows on
measures, clamped scales and per-volume constants, and depth-zero
circularity --- and with the typed repairs that exclude each.  The
central hypothesis, \lean{PhysicalTerminalScaleWilsonGate}, is an open
proposition: no witness is provided, and the infrared/ultraviolet
bounds it demands of the actual Wilson measure are exactly the open
analytic mathematics (Balaban-type single-scale estimates).  The final
theorem is conditional: a witness of the gate yields
$|\Cov(2^n u)| \le C e^{-m \cdot 2^n u}$ with one pair $(C, m)$,
$m > 0$, for every base $M_0 \ge 4$ and every depth $n \ge 1$.  No
mass-gap claim, no claim of gate satisfiability, and no thermodynamic
or continuum limit is made or implied.
\end{abstract}

\section{The specification problem}\label{sec:spec}

The renormalization-group analysis of lattice gauge theories in the
tradition of Balaban \cite{Bal87,Bal88a,Bal88b,Bal89} (see Dimock
\cite{Dim13a,Dim13b,Dim14} for a modern exposition, and
\cite{Cha19,Cha26} for the probabilistic state of the art) organizes a
correlation function of the fundamental measure into an infrared part,
carried by an effective theory at a coarse terminal scale, plus a sum
of per-scale ultraviolet remainders.  This paper concerns a question
that is logically prior to any estimate: \emph{what does it mean, in a
proof assistant, to state that organization faithfully?}

The question is not vacuous.  A formalized ``multiscale
decomposition'' of a correlator function
$t \mapsto \mathrm{cov}(t)$ of the shape
\begin{equation}\label{eq:decomp}
\mathrm{cov}(t) \;=\; \mathrm{cov}_{\mathrm{IR}}(t)
  \;+\; \sum_{k < n_{\mathrm{sc}}(t)} R_{t,k},
\end{equation}
with exponential bounds on $\mathrm{cov}_{\mathrm{IR}}$ and per-scale
bounds on the $R_{t,k}$, is \emph{gameable} whenever the decomposition
data $(\mathrm{cov}_{\mathrm{IR}}, R, n_{\mathrm{sc}})$ is quantified
freely: known information can be relabeled into the required slots
with zero renormalization-group content, and the resulting
``theorems'' transport hypotheses to conclusions without ever touching
the physics.  A second, independent trap is volumetric: on a fixed
finite torus a two-plaquette correlator realizes only finitely many
separations, so \emph{every} such function satisfies \emph{some}
exponential bound at a large enough constant; interfaces that
quantify their constants per volume
($\forall N\, \exists C(N)\, \dots$) are therefore contentless in the
direction of the thermodynamic regime, and the quantifier order
$\exists C\, \forall N \dots$ must be enforced \emph{in the type} of
the interface.

The contribution of this paper is a mechanized interface for the
concrete $\SU(N_c)$ Wilson theory that survives its own adversarial
audit.  The development method was to state each candidate interface,
attack it, record the explicit countermodel when the attack succeeded,
and repair the interface \emph{in the type} so that the countermodel
becomes unstatable or provably excluded --- then iterate.
Section~\ref{sec:counter} presents the four rounds of this process as
mathematics: each naive interface, its countermodel with the explicit
formula, and the typed repair.  Sections~\ref{sec:rg}--\ref{sec:oper}
present the surviving architecture: the concrete decimation and its
proved properties (\S\ref{sec:rg}), the non-circular terminal gate
with the physical-distance exponent (\S\ref{sec:gate}), the transported
support geometry (\S\ref{sec:support}), and the operational
certification composing them (\S\ref{sec:oper}).
Section~\ref{sec:main} states the final conditional theorem;
Section~\ref{sec:frontier} names what remains open, plainly;
Section~\ref{sec:repro} records the reproducibility data and the
theorem-to-artifact map.

\medskip
\noindent\textbf{Statuses.}  As in the companion papers of this
repository \cite{ErikssonC2,ErikssonC3,ErikssonC4,ErikssonC5}, we use
a strict status discipline: \textbf{exact} means a Lean~4 theorem in
the pinned central repository whose axiom oracle prints exactly
$[\lean{propext}, \lean{Classical.choice}, \lean{Quot.sound}]$;
\textbf{paper-level} means ordinary prose, always flagged.  Everything
in Sections~\ref{sec:rg}--\ref{sec:main} is exact unless flagged; the
few paper-level items (one geometric counterexample in
\S\ref{sec:support}) support no theorem.  The development contains no
\lean{sorry} and no project axioms; open questions are carried as
named \emph{open propositions} --- definitions of \lean{Prop}s with no
witness --- and are inventoried in \S\ref{sec:frontier} and
Table~\ref{tab:open}.

\medskip
\noindent\textbf{Honest scope, stated at the outset.}  All objects are
finite-torus objects.  The interface's central hypothesis has
\emph{no witness in this development}; providing one for the Wilson
data is exactly the open analytic problem (the Balaban-type
single-scale activity estimate), and nothing below depends on its
satisfiability.  The conditional conclusion is a mass-gap-\emph{shaped}
decay bound on a finite-torus correlator family; no thermodynamic or
continuum limit is constructed, and no claim on the Clay Yang--Mills
problem is made (\S\ref{sec:frontier}).  On the formalization side,
the closest antecedent we know is the Lean~4 formalization of the
\emph{free} scalar field and its Glimm--Jaffe axioms by Douglas,
Hoback, Mei and Nissim \cite{DHMN26}; the present development is
complementary in ambition and scope --- it formalizes an
\emph{interacting-theory renormalization-group interface} (definitions,
identities, geometry, and conditional transport for Wilson lattice
gauge theory), not an estimate and not a quantum field theory.  Other
formalized material in the same repository includes a volume-uniform
Wilson-loop area law via a formalized cluster expansion
\cite{ErikssonAL}, rooted-tree majorants for polymer expansions with
holes \cite{ErikssonRT}, and exact two-dimensional $\SU(2)$
Yang--Mills computations \cite{Eriksson2D}; the strong-coupling
clustering input quoted in \S\ref{sec:main} belongs to that cluster
expansion lane.

\section{Naive interfaces and countermodels}\label{sec:counter}

Throughout, $\mathrm{cov} : \mathbb{N} \to \mathbb{R}$ is a correlator
function --- in the intended instance, a distance-indexed truncated
two-plaquette correlator of the Wilson Gibbs measure
(\S\ref{sec:rg:setting}).  The four interfaces below are presented in
the order of their design; each is retained in the development (with
corrected documentation) because the repaired interfaces factor
through them, but \emph{none of the first three is load-bearing}: they
are the audit trail.  All countermodel formulas below are recorded in
the development's documentation; the clamp and depth-zero analyses of
\S\ref{sec:cm3}--\ref{sec:cm4} are additionally \emph{formalized} as
the exact theorems cited there.

\subsection{The abstract decomposition, and scalar relabeling}
\label{sec:cm1}

\begin{interface}[abstract bridge; \lean{CorrelatorBridge},
\lean{RegimeCoherentWilsonBridge}]\label{int:abstract}
A \emph{bridge} for $\mathrm{cov}$ is a tuple
$(\mathrm{cov}_{\mathrm{IR}}, R, n_{\mathrm{sc}})$ satisfying the
identity \eqref{eq:decomp} for all $t$.  The naive gate asks: there
exists a bridge, a marginal-coupling flow $g$
(Definition~\ref{def:flow} below), and constants such that
$|\mathrm{cov}_{\mathrm{IR}}(t)| \le C_1 e^{-\varepsilon t}$, each
$|R_{t,k}| \le C_2 e^{-c_0 t} g_k^{\kappa_0}$, and $R \not\equiv 0$.
\end{interface}

The transport theorem for this interface
(\lean{physical\_mass\_gap\_of\_bridge}) is sound: any witness yields
$|\mathrm{cov}(t)| \le (C_1 + C_2 \sum_k g_k^{\kappa_0})
e^{-\min(\varepsilon, c_0)\, t}$.  The interface is nevertheless
worthless as a statement of RG content:

\begin{countermodel}[scalar relabeling]\label{cm:relabel}
Suppose some decay $|\mathrm{cov}(t)| \le a e^{-t}$ is already known
(for the Wilson theory at strong coupling, such a bound is a theorem
--- see \S\ref{sec:main}).  Set
\[
R_{t,0} := a e^{-t}, \qquad R_{t,k} := 0 \ (k \ge 1), \qquad
\mathrm{cov}_{\mathrm{IR}}(t) := \mathrm{cov}(t) - a e^{-t}, \qquad
n_{\mathrm{sc}} := 1 .
\]
Every clause of Interface~\ref{int:abstract} is satisfied --- the
decomposition identity trivially, the IR bound by the triangle
inequality, the UV bound after absorbing $g_0^{-\kappa_0}$ into $C_2$,
and the nonvanishing clause by $R_{t,0} \ne 0$ --- with \emph{zero}
renormalization-group content.  The clause ``$R \not\equiv 0$''
certifies nonvanishing, not ultraviolet provenance.
\end{countermodel}

\begin{remark}[typed repair: remainders must be produced, not selected]
\label{rem:repair1}
The repair removes the freedom to choose $R$: introduce a family of
measures $\mu_0, \mu_1, \mu_2, \dots$ on gauge configurations with
$\mu_0$ the Wilson Gibbs measure, define the \emph{effective
correlator} $\mathrm{cov}_k$ at scale $k$ as the correlator \emph{of}
$\mu_k$ --- a definition, not data --- and \emph{produce}
\[
R_{t,k} := \mathrm{cov}_k(t) - \mathrm{cov}_{k+1}(t), \qquad
\mathrm{cov}_{\mathrm{IR}}(t) := \mathrm{cov}_{n_{\mathrm{sc}}(t)}(t).
\]
The decomposition \eqref{eq:decomp} is then a \emph{theorem}
(telescoping; \lean{provenant\_decomposition}), not a field, and a
nonzero remainder certifies that the measure genuinely changed between
consecutive scales
(\lean{step\_moves\_of\_provenantRsc\_ne\_zero}).
Artifacts: \lean{RGFlowFamily}, \lean{covEff}, \lean{provenantRsc},
\lean{provenantBridge}.
\end{remark}

\subsection{Measure-level flows, and the sink flow}\label{sec:cm2}

Remark~\ref{rem:repair1} still quantifies over the \emph{generating
map} of the family: an \lean{RGFlowFamily} carries a field
$T : \mathrm{Measures} \to \mathrm{Measures}$ with
$\mu_{k+1} = T(\mu_k)$.  That field is the next attack surface.

\begin{countermodel}[sink flow]\label{cm:sink}
Take $T(\nu) := 0$ for every $\nu$ (or a fixed Dirac measure
$\delta_{A_0}$, whose connected correlators vanish, so that
probability preservation alone would not repair the interface).  Then
$\mathrm{cov}_0 = \mathrm{cov}$ and $\mathrm{cov}_k = 0$ for
$k \ge 1$: the \emph{entire} physical correlator is relabeled as the
first-scale ultraviolet remainder $R_{t,0} = \mathrm{cov}(t)$, its
known strong-coupling bound converts into the demanded single-scale
shape by adjusting $C_2$ by $g_0^{-\kappa_0}$, and the terminal
infrared part is identically zero.  This is
Countermodel~\ref{cm:relabel} one level up: an arbitrary total map of
measures can encode \emph{any} prescribed orbit.
\end{countermodel}

\begin{remark}[typed repair: no map variable]\label{rem:repair2}
The repair removes the map variable altogether.  The scale
transformation is a \emph{fixed, defined} decimation
$B_M$ acting on gauge \emph{configurations}
(Definition~\ref{def:blockmap}), and the effective measure is its
literal pushforward, $\mathcal{E}_M \mu := (B_M)_* \mu$ --- by
construction, not as data.  Nothing in the surviving interfaces
quantifies over transformations, so the sink flow is
\emph{unstatable}; and probability preservation is a theorem
(\lean{effectiveMeasure\_isProbability}), so no stage of any tower can
be the zero measure (\lean{effectiveMeasure\_ne\_zero}).  The fixed
Dirac is equally unreachable: the only measures occurring in the
surviving interfaces are iterated pushforwards of the Wilson Gibbs
measure itself.  If the honestly blocked measures happened to have
vanishing correlator differences, the nonvanishing data clause of the
gate would simply be \emph{false} --- the attack can prevent a witness,
never fabricate one.
\end{remark}

\subsection{Clamped scales and per-volume constants}\label{sec:cm3}

The decimation halves the torus, so a multiscale tower must be indexed
carefully: starting from size $2^n M_0$, after $n$ steps it reaches
the base size $M_0$ and can descend no further.  The development makes
the tower a \emph{total} function by clamping --- beyond depth $n$ the
tower stays at its base stage (\S\ref{sec:rg:tower}).  Totality is a
formalization convenience; unguarded, it is an attack surface, and it
interacts with the volumetric trap of \S\ref{sec:spec}.

\begin{countermodel}[clamp and per-volume constants]\label{cm:clamp}
(i) In a gate whose scale count $n_{\mathrm{sc}}$ is free data, a
witness may point the ``terminal'' infrared object at a clamped index
($n_{\mathrm{sc}}(t) > n$, where the tower is stationary) or at the
root ($n_{\mathrm{sc}}(t) = 0$, where no blocking has occurred and
the ``infrared part'' \emph{is} the physical correlator): relabeling
reappears through the index arithmetic.  Similarly, a depth-$0$ tower
($n = 0$) produces no remainders at all, so ``nontriviality'' clauses
quantified over all depths are vacuously threatened at $n = 0$.
(ii) Independently, on the fixed torus of size $N$ the canonical
correlator (\S\ref{sec:rg:setting}) is supported on finitely many
separations, so for \emph{each} $N$ some constant $C(N)$ makes any
exponential bound true; a gate of quantifier shape
$\forall N\, \exists C(N)$ is therefore satisfiable without content.
\end{countermodel}

\begin{remark}[typed repair: clamp theorems and quantifier order]
\label{rem:repair3}
Against (i), the clamp is \emph{defused by theorems}: clamped scales
produce \emph{zero} remainder
(\lean{concreteRsc\_eq\_zero\_of\_le},
\lean{scaledCanonicalRsc\_eq\_zero\_of\_le}), so a nonzero produced
remainder forces $k < n$ --- ultraviolet content certifies an actual
decimation step at an actual genuine scale
(\lean{lt\_of\_concreteRsc\_ne\_zero},
\lean{lt\_of\_scaledCanonicalRsc\_ne\_zero}); nonvanishing clauses are
restricted to $n \ge 1$; and the terminal index is ultimately
\emph{pinned by definition} with a typed range
(\S\ref{sec:gate}), so the infrared clause cannot reach the clamp
region or the root at all.  Against (ii), every surviving gate has the
quantifier shape
$\exists\, (\text{constants})\ \forall\, (\text{base } M_0)\
\forall\, (\text{depth } n) \dots$ --- one constant pack before the
unbounded volume family (the tower sizes are strictly increasing and
cofinal in all sizes: \lean{towerSize\_strictMono},
\lean{towerSize\_unbounded}); the forbidden shape is unstatable
against the gate's type, and the consumer's conclusion
(Theorem~\ref{thm:main}) carries the same order.
\end{remark}

\subsection{Depth-zero circularity}\label{sec:cm4}

One further failure mode survives all of the above, and it is the
subtlest: a gate can be non-gameable and still contain its intended
conclusion \emph{as a hypothesis}, if its clause family includes a
layer at which the tower performs no blocking.

\begin{countermodel}[the $n = 0$ shell]\label{cm:circular}
Consider the ``literal'' gate that demands the scale-corrected bounds
at every base $M_0 \ge 4$ \emph{and every depth $n \ge 0$}
(\lean{LiteralFidelityConcreteRGWilsonGate}).  Every torus size
$N \ge 4$ occurs in this family as a base at depth zero
($M_0 := N$, $n := 0$), and at depth zero the scaled infrared clause
reads: \emph{the physical correlator itself decays exponentially at
every separation in the window, uniformly in $N$} --- which is the
statement the whole architecture is trying to reach.  Consequently
the implication ``literal gate $\Rightarrow$ decay at all separations
and all sizes'' (\lean{allSeparations\_of\_literalFidelityGate}, a
theorem) consumes \emph{only} the depth-zero shell and no
renormalization-group content whatsoever.  The implication is valid;
as a reduction of the open problem it is circular.  Two exact
factorization facts sharpen the diagnosis: an odd separation $t$ lies
on the depth-$n$ dyadic separation lattice $\{2^n u\}$ only at
$n = 0$ (\lean{odd\_dyadic\_factorization\_trivial},
\lean{odd\_separation\_needs\_base\_scale}), and no separation
$t < 2^n$ is reachable at depth $n$ at all
(\lean{dyadic\_lattice\_misses\_small\_separations}); so a gate that
controls all separations \emph{must} lean on its depth-zero layer ---
there is no honest way to reach odd $t$ from positive scales without
new interpolation input (\S\ref{sec:frontier}).
\end{countermodel}

\begin{remark}[typed repair: the fixed terminal index]
\label{rem:repair4}
The repair excludes the depth-zero layer in the type.  The terminal
gate of \S\ref{sec:gate} quantifies \emph{only} depths $n \ge 1$; its
infrared clause is evaluated \emph{syntactically} at the fixed
terminal index $\kT(n) := n$, with the typed range
$1 \le \kT(n) \le n$ (\lean{kTerm\_pos}, \lean{kTerm\_le} --- both
definitional), so the consumed infrared object is the correlator of
the \emph{maximally blocked} measure, never of the unblocked one; and
the index $k = 0$ enters the gate only \emph{inside the first
difference} $R_0(u) = \mathrm{cov}_0(u) - \mathrm{cov}_1(u)$
(\lean{rsc\_is\_difference}, by \lean{rfl}) --- a single-scale
remainder of the honest ultraviolet shape, not a standalone decay
hypothesis on the physical object.  The bound on the physical
correlator is then \emph{derived} through the telescoping, not
hypothesized (Theorem~\ref{thm:main}).
\end{remark}

\section{The concrete gauge renormalization group}\label{sec:rg}

\subsection{Setting}\label{sec:rg:setting}

Fix a dimension $d \ge 2$.  For $N \ge 1$ let
$\Torus{N} = (\mathbb{Z}/N\mathbb{Z})^d$ be the periodic torus, with
oriented edges $(y, i, \pm)$ ($y \in \Torus{N}$, $i \in \{0, \dots,
d-1\}$).  For a group $G$ (measurable; the physical instance is the
compact $\SU(N_c)$ with its Borel structure), a \emph{gauge
configuration} is a $G$-valued function $A$ on oriented edges obeying
the reversal constraint $A(e^{-1}) = A(e)^{-1}$; the configuration
space is $\Conf{N} = \Conf{N}(G)$, with the product measurable
structure of the positive-edge coordinates.  A \emph{plaquette}
$p = (y, i, j)$, $i < j$, has the usual four-edge boundary and
holonomy $\holo(A, p) \in G$; gauge transformations
$u : \Torus{N} \to G$ act by
$(u \cdot A)(y, i, +) = u(y)\, A(y, i, +)\, u(y + e_i)^{-1}$.

The \emph{Wilson Gibbs measure} at coupling $\beta \in \mathbb{R}$ is
the tilted product Haar measure
\begin{equation}\label{eq:wilson}
d\muW_{N,\beta}(A) \;=\;
Z_{N,\beta}^{-1}\,
e^{-\beta \sum_{p} E(\holo(A,p))}\; d\mathrm{Haar}^{\otimes}(A),
\qquad E(U) = \mathrm{Re}\,\mathrm{tr}(\mathbf{1}_{N_c} - U),
\end{equation}
realized as an honest \lean{Measure.tilted}
(\lean{wilsonGibbsMeasure}); it is a probability measure for every
torus size and every real $\beta$
(\lean{wilsonGibbsMeasure\_isProbability}), and its correlators agree
exactly, with no hypotheses, with the $Z$-quotient Boltzmann
expressions (\lean{integral\_gibbsMeasure\_eq\_div},
\lean{rayCorrelatorOfMeasure\_wilsonGibbs}).

For a measure $\nu$ on $\Conf{N}$ and an observable
$f : G \to \mathbb{R}$, the \emph{truncated two-plaquette correlator}
is
\[
T_\nu(p, q) \;=\; \int f(\holo(A,p))\, f(\holo(A,q))\, d\nu
 \;-\; \int f(\holo(A,p))\, d\nu \int f(\holo(A,q))\, d\nu .
\]
Plaquettes carry the \emph{touch graph}: $p \sim q$ iff $p \ne q$ and
their supports share an edge; $\dist_N$ denotes its graph distance.
The \emph{canonical plaquette family} kills all selection freedom: for
$t \in \mathbb{N}$, $P^N_t$ is the axis-aligned plaquette in the
$(e_0, e_1)$-plane based at the site $t \cdot e_0$ (torus wrap by
$t \bmod N$) --- a closed-form geometric definition
(\lean{canonicalPlaquette}), with exact linear metric growth
\begin{equation}\label{eq:candist}
\dist_N(P^N_0, P^N_t) = t \qquad (2t \le N)
\end{equation}
(\lean{canonicalPlaquette\_dist\_eq}; the lower bound is proved by a
$\mathbb{Z}/N$-valued potential argument, wrap-safe, the upper bound
by an explicit axis walk).  The \emph{canonical correlator} at
separation index $t$ is
\[
\Cov_\nu(t) \;:=\; T_\nu\bigl(P^N_0,\; P^N_{2t}\bigr),
\]
so that in the wrap-free window $1 \le t$, $4t \le N$ the pair is
distinct at touch-distance exactly $2t$
(\lean{canonicalPair\_admissible}).  (The development also carries a
choice-based ``ray'' correlator and an all-pairs supremum correlator
\lean{supAbsCorrelator}, with domination lemmas connecting the three;
the canonical family is the load-bearing one.)

\subsection{The decimation}\label{sec:rg:blockmap}

\begin{definition}[the block map; \lean{blockMap}]\label{def:blockmap}
For $M \ge 1$, the decimation
$B_M : \Conf{2M}(G) \to \Conf{M}(G)$ sends a fine configuration $A$ to
the coarse configuration whose value on the positive coarse edge
$(y, i, +)$ is the ordered product of the two fine holonomies along
its doubled image:
\[
(B_M A)(y, i, +) \;=\; A(2y, i, +) \cdot A(2y + e_i,\, i, +),
\]
negative edges receiving the inverse through the reversal constraint.
Here $y \mapsto 2y$ is the even-sublattice embedding
(\lean{coarseSiteEmbed}).
\end{definition}

$B_M$ is a fixed definition --- there is no transformation variable
anywhere downstream (Remark~\ref{rem:repair2}).  Four properties are
proved about it as theorems:

\begin{proposition}[properties of the decimation]\label{prop:blockmap}
\begin{enumerate}
\item[(i)] \emph{Measurability}: the map $B_M$ is measurable when
multiplication on $G$ is jointly measurable
(\lean{measurable\_blockMap}).
\item[(ii)] \emph{Locality} (\lean{blockMap\_local}): the coarse value
at an edge depends only on the fine configuration at its two
constituent fine edges --- finite range, with the explicit two-edge
carrier named in the hypotheses.
\item[(iii)] \emph{Gauge covariance} (\lean{blockMap\_gaugeAct}):
$B_M(u \cdot A) = (u \circ \mathrm{embed}) \cdot B_M(A)$ --- blocking a
gauge-transformed configuration equals gauge-transforming the blocked
configuration by the restriction of $u$ to the even sublattice (the
interior vertex of each doubled edge cancels).
\item[(iv)] \emph{Scale geometry} (\lean{coarseSiteEmbed\_shift}): one
coarse step equals two fine steps, torus wrap included.
\end{enumerate}
\end{proposition}

\begin{definition}[effective measure; \lean{effectiveMeasure}]
$\mathcal{E}_M \mu := (B_M)_* \mu$, the pushforward.  It preserves
probability (\lean{effectiveMeasure\_isProbability}) and in particular
is never the zero measure (\lean{effectiveMeasure\_ne\_zero}).
\end{definition}

Observable transport is proved in exactly the correlator shape the
interface consumes: the coarse integral, truncated correlator, and
distance-indexed correlator under $\mathcal{E}_M \mu$ equal the fine
integrals of the pulled-back observable
$A \mapsto f(\holo(B_M A, p))$ under $\mu$
(\lean{integral\_effectiveMeasure},
\lean{truncatedCorrelator\_effectiveMeasure},
\lean{rayCorrelator\_effectiveMeasure}), via the change-of-variables
identity for pushforwards with the measurability side conditions
discharged from Proposition~\ref{prop:blockmap}(i).

\subsection{The tower and the scale-corrected telescoping}
\label{sec:rg:tower}

Fix a base size $M_0 \ge 1$ and define tower sizes by the doubling
recursion $\mathrm{sz}(0) = M_0$,
$\mathrm{sz}(m+1) = 2\, \mathrm{sz}(m)$, so
$\mathrm{sz}(m) = 2^m M_0$ (\lean{towerSize},
\lean{towerSize\_eq\_pow\_mul}); the recursion (rather than the closed
form) makes the dependently-indexed tower typecheck without casts.
For a start measure $\mu$ on $\Conf{2^n M_0}$, the \emph{tower
measure} $\mu^{(k)}$ is the $k$-fold blocked measure, living on
$\Conf{2^{n-k} M_0}$, defined by iterating $\mathcal{E}$; at the base
scale the tower \emph{clamps} (stays put) --- stated openly and
defused by the clamp theorems of Remark~\ref{rem:repair3}.  Every
stage of a Wilson-anchored tower is a typed probability measure
(\lean{towerMeasure\_isProbability},
\lean{wilsonTowerMeasure\_isProbability}).

Stage $k$ of the tower lives on a lattice whose spacing is $2^k$ in
finest units.  The \emph{scale-corrected canonical correlators}
therefore evaluate stage $k$ at separation index $2^{n-k} u$:
\begin{equation}\label{eq:scaled}
c_k(u) \;:=\; \Cov_{\mu^{(k)}}\bigl(2^{n-k} u\bigr), \qquad
R_k(u) \;:=\; c_k(u) - c_{k+1}(u),
\end{equation}
(\lean{scaledCanonicalCovEff}, \lean{scaledCanonicalRsc}), so that
every stage sees the \emph{same} physical separation index
$2^k \cdot (2^{n-k} u) = 2^n u$, and the separation-to-volume ratio
$u / M_0$ is scale-invariant: the single window
\begin{equation}\label{eq:window}
1 \le u, \qquad 4u \le M_0
\end{equation}
serves the whole tower, with the canonical pair realized at
touch-distance exactly twice the stage separation index at every stage
(\lean{scaledCanonical\_pair\_dist}).  The multiscale decomposition is
then a theorem:

\begin{proposition}[telescoping; \lean{scaledCanonical\_decomposition}]
\label{prop:telescope}
For every scale count $n_{\mathrm{sc}}$ and every $u$,
\[
c_0(u) \;=\; c_{n_{\mathrm{sc}}(u)}(u)
  \;+\; \sum_{k < n_{\mathrm{sc}}(u)} R_k(u),
\]
and $c_0(u) = \Cov_{\mu}(2^n u)$, the canonical correlator of the
\emph{base} measure at the physical separation index
(\lean{scaledCovEff\_zero\_physical}, definitional).
\end{proposition}

Nothing in \eqref{eq:scaled} is selectable: the only inputs are the
Wilson data $(\mu, f)$ and the tower geometry $(M_0, n)$.

\section{The non-circular terminal gate}\label{sec:gate}

\begin{definition}[marginal-coupling flow]\label{def:flow}
A \emph{marginal flow} is a sequence $g : \mathbb{N} \to (0, \infty)$
together with $\beta_c > 0$ such that $\beta_c g_k < 1$ and
$g_{k+1} = g_k (1 - \beta_c g_k)$ for all $k$ --- the discrete
logistic recursion mimicking an asymptotically free marginal coupling.
For every $\kappa_0 > 1$ the sum $\sum_k g_k^{\kappa_0}$ converges;
this is a theorem of the development
(\lean{marginal\_coupling\_pow\_summable\_of\_recursion}), and it is
what turns per-scale bounds into a summed constant.
\end{definition}

\begin{definition}[terminal index; \lean{kTerm}]\label{def:kterm}
$\kT(n) := n$.  The range facts $1 \le \kT(n)$ (for $n \ge 1$) and
$\kT(n) \le n$ are definitional (\lean{kTerm\_pos}, \lean{kTerm\_le}):
the terminal stage is the maximally blocked measure, the clamp region
is unreachable, and the root $k = 0$ is not the infrared object
(Remark~\ref{rem:repair4}).
\end{definition}

\begin{definition}[the physical terminal gate;
\lean{PhysicalTerminalScaleWilsonGate}]\label{def:gate}
Fix $d \ge 2$, $N_c \ge 1$, an observable
$f : \SU(N_c) \to \mathbb{R}$, and $\beta \in \mathbb{R}$.  The gate
holds iff there exist a marginal flow $(g, \beta_c)$ and constants
$C_1 \ge 0$, $C_2 \ge 0$, $\varepsilon > 0$, $c_0 > 0$,
$\kappa_0 > 1$ --- \emph{all quantified before every base and depth}
--- such that for every base $M_0 \ge 4$ and every depth $n \ge 1$,
writing $\mu = \muW_{2^n M_0, \beta}$ and $c_k, R_k$ for the
scale-corrected canonical tower correlators \eqref{eq:scaled} of the
Wilson-anchored tower:
\begin{itemize}
\item[(P)] every tower stage is a probability measure --- a clause
that is unconditionally provable
(\lean{fidelityGate\_probability\_clause}) and is carried typed so
that every admitted instance is a probability tower by its statement;
\item[(IR)] for every $u$ in the window \eqref{eq:window},
\[
\bigl| c_{\kT(n)}(u) \bigr|
 \;\le\; C_1\, e^{-\varepsilon \cdot 2^n u} ,
\]
where the left-hand side is \emph{syntactically} the terminal object
at the index $\kT(n)$ (\lean{terminalGate\_IR\_index\_eq}, by
\lean{rfl}: same term, no bridging freedom);
\item[(UV)] for every such $u$ and every $k < \kT(n)$,
\[
\bigl| R_k(u) \bigr|
 \;\le\; C_2\, e^{-c_0 \cdot 2^n u}\, g_k^{\kappa_0} ;
\]
\item[(NZ)] some $R_k(u) \ne 0$ with $u$ in the window and
$k < \kT(n)$ --- a data clause for the given Wilson data, not
engineered into the class (a genuine RG may have fixed measures and
invariant observables; for constant $f$ all correlators vanish and
(NZ) is simply false).
\end{itemize}
\end{definition}

Three design points, each enforced in the type.  First, the
\emph{decay variable is physical}: the exponents evaluate at
$2^n u$ --- the cast $((2^n u : \mathbb{N}) : \mathbb{R})$ appears
literally in the Lean \lean{Prop} --- so a witness must provide a gap
that does not shrink with the depth; the variant gate with
renormalized exponents $e^{-\varepsilon u}$
(\lean{TerminalScaleWilsonGate}) is strictly weaker, and the
implication physical $\Rightarrow$ renormalized is a theorem
(\lean{terminalScaleGate\_of\_physicalTerminalGate}) whose converse is
false pointwise and is neither stated nor consumed.  Second, the
quantifier order is the one demanded by
Countermodel~\ref{cm:clamp}(ii).  Third, the window \eqref{eq:window}
is wrap-free and depth-independent: it is a condition on $(u, M_0)$
only (\lean{terminal\_window\_nonempty}), so it does not shrink as $n$
grows.

\medskip
\noindent\textbf{No witness.}  No witness of
Definition~\ref{def:gate} is provided anywhere in this development,
and nothing below depends on its satisfiability.  Producing the
(IR)/(UV) clauses for the actual Wilson measure is the open analytic
problem --- in substance, a Balaban-type single-scale estimate for the
blocked Wilson theory \cite{Bal87,Bal88a,Bal88b,Bal89,Dim13a}; the
gate is the typed statement of \emph{what exactly} must be proved, in
a form that Sections~\ref{sec:counter} and \ref{sec:oper} certify
cannot be discharged by relabeling.

\section{Support geometry}\label{sec:support}

The telescoping of Proposition~\ref{prop:telescope} is honest algebra,
but a scale decomposition of \emph{one physical separation} requires
metric control: the selected plaquette pairs at consecutive scales
must be the same geometric object at matched offsets, and the
transported observables must have supports that genuinely separate.
This section records the proved geometry.

\subsection{One step: embedding and the exact canonical doubling}

The plaquette embedding
$\iota : p \mapsto (\text{site } 2y, \text{ same plane})$
(\lean{plaqEmbed}) is injective.  The naive expectation
$\dist_{2M}(\iota p, \iota q) \le 2 \dist_M(p, q)$ is \emph{false}:
in $d = 3$ there are coarse plaquettes at touch-distance $1$ (sharing
one edge, in transverse planes) whose embedded images sit at fine
touch-distance exactly $3$ --- a single fine plaquette cannot meet
both embedded supports, while an explicit length-$3$ walk through the
doubled shared edge exists.  (This counterexample is a paper-level
analysis recorded in the development's documentation; no theorem
depends on it --- it explains why the general theorem has the constant
it has.)  What is proved:

\begin{proposition}[metric transport, one step]\label{prop:onestep}
\begin{enumerate}
\item[(i)] For touch-reachable pairs,
$\dist_{2M}(\iota p, \iota q) \le 3\, \dist_M(p, q)$
(\lean{plaqEmbed\_dist\_le\_three\_mul}), with walks transported
constructively.  No sharpness claim is formalized.  The converse
direction (coarse distance dominated by embedded fine distance) is a
named open proposition, unproved and consumed nowhere
(\lean{PlaqEmbedDistLowerBound}).
\item[(ii)] On the canonical family the scale relation is exact with
the physical factor: $\iota(P^M_\tau) = P^{2M}_{2\tau}$
(\lean{plaqEmbed\_canonicalPlaquette}) and, for $2\tau \le M$,
\[
\dist_{2M}\bigl(\iota P^M_0, \iota P^M_\tau\bigr)
 \;=\; 2\, \dist_M\bigl(P^M_0, P^M_\tau\bigr)
\]
(\lean{canonicalPlaquette\_dist\_doubles}).  The interface consumes
only this family.
\end{enumerate}
\end{proposition}

\subsection{One step: the support of the pullback observable}

The pullback observable $A \mapsto f(\holo(B_M A, p))$ reads the fine
configuration only on the eight fine edges of the $2 \times 2$ fine
loop over $p$ --- the set $\mathrm{supp}_B(p)$
(\lean{blockPlaquetteSupport}, $|\mathrm{supp}_B(p)| \le 8$), with the
locality statement proved in the observable shape the correlators
consume (\lean{blockMap\_holonomy\_congr},
\lean{blockObservable\_congr}).  Every support edge's source lies in
the $2 \times 2$ plane box at the doubled anchor
(\lean{blockPlaquetteSupport\_source\_in\_ball}), and for the
canonical pair at coarse offset $\tau$ in the window
$1 \le \tau$, $4\tau \le M$, \emph{any} fine plaquettes whose supports
meet the two transported supports are separated by at least
$2\tau - 3$ touch-steps
(\lean{blockSupport\_canonical\_separation\_walk},
\lean{blockSupport\_canonical\_separation\_dist}) --- the slack $3$
being the radius-$2$ support extent plus one carrier-internal offset,
proved by the same wrap-safe circular-potential argument as
\eqref{eq:candist}.

\subsection{$k$ steps: radius, separation, positivity}

Iterating, the $k$-step transported support of a plaquette
(\lean{kStepBlockSupport}) is controlled by a master induction:

\begin{theorem}[$k$-step radius;
\lean{kStepBlockSupport\_source\_in\_ball}]\label{thm:radius}
Every edge of the $k$-step transported support of a plaquette $P$
keeps its direction in the plane $\{ \mathrm{dir}_1 P,
\mathrm{dir}_2 P\}$, and its source lies in the plane box at the
$2^k$-scaled anchor (\lean{iterEmbed\_coord}: the anchor coordinates
scale exactly, no wrap) with direction-correlated offsets
$a + [\mathrm{dir} = \mathrm{dir}_1] \le 2^k$,
$b + [\mathrm{dir} = \mathrm{dir}_2] \le 2^k$.  In particular the
support radius is at most $2^k$, and one step of the recursion is
definitionally the one-step object
(\lean{kStepBlockSupport\_one}).
\end{theorem}

\begin{theorem}[$k$-step separation;
\lean{kStepSupport\_canonical\_separation\_walk},
\lean{kStepSupport\_canonical\_separation\_dist}]\label{thm:ksep}
Let $1 \le \tau$ and $2\tau \le M$ with $M$ a tower size.  Any fine
plaquettes on the size-$2^k M$ torus whose supports meet the two
$k$-step transported supports of the canonical pair
$(P^M_0, P^M_\tau)$ are separated by at least
\[
2^k \tau - (2^k + 1)
\]
touch-steps, on every touch-walk (hence in graph distance when
reachable).  At $k = 1$ this is the one-step constant $2\tau - 3$
exactly (\lean{kStepSeparation\_one}).  On the window $\tau \ge 2$,
$k \ge 1$ the bound is strictly positive
(\lean{kStepSeparation\_pos}) and strictly increasing in $k$
(\lean{kStepSeparation\_strictMono}); at $\tau = 1$ it
$\mathbb{N}$-truncates to $0$ (true and uninformative, stated), and
neither positivity nor sharpness is claimed outside the stated window.
\end{theorem}

\subsection{The terminal instance}

The gate of Definition~\ref{def:gate} consumes the terminal stage
$k = \kT(n) = n$, whose infrared object is --- definitionally, by
\lean{rfl} --- the truncated correlator of the $n$-fold blocked
measure at the canonical pair $(P^{M_0}_0, P^{M_0}_{2u})$ on the
\emph{base} torus (\lean{terminalIR\_is\_canonical\_pair\_correlator},
with \lean{terminal\_stage\_size} and
\lean{terminal\_stage\_sep\_index}).  Instantiating
Theorem~\ref{thm:ksep} at exactly that pair and exactly $k = n$
(coarse offset $\tau = 2u \ge 2$) gives: in the gate's own window
\eqref{eq:window}, the $n$-step transported supports of the terminal
pair separate every pair of touching carriers by at least
\begin{equation}\label{eq:termsep}
2^n (2u) - (2^n + 1) \;>\; 0 \qquad (n \ge 1,\ u \ge 1),
\end{equation}
with strict growth in the depth
(\lean{terminalSupport\_canonical\_separation\_walk},
\lean{terminalSupport\_separation\_pos},
\lean{terminalSupport\_separation\_strict\_growth}).  The window is a
condition on $(u, M_0)$ only --- no $n$ appears in the hypotheses.

\section{Operational certification}\label{sec:oper}

Section~\ref{sec:support} is geometry of \emph{sets};
Definition~\ref{def:gate} consumes \emph{correlators of measures}.
The remaining obligation --- and the last piece of the architecture
--- is the operational connection: that the certified geometry is the
geometry \emph{of the observable the gate consumes}.  This is supplied
by one composed map and three theorems about it.

\begin{definition}[the composed block map; \lean{iteratedBlockMap}]
$B^{(k)}_{M_0, n} : \Conf{2^n M_0} \to \Conf{2^{n-k} M_0}$ is the
$k$-fold composition of the tower's configuration-level steps, by the
\emph{same} recursion as the tower measure (one definition; every
theorem below consumes this term).  One step on a genuine scale is
definitionally $B_M$ (\lean{iteratedBlockMap\_one}); the map is
measurable (\lean{measurable\_iteratedBlockMap}).
\end{definition}

\begin{theorem}[measure identity;
\lean{towerMeasure\_eq\_map\_iterated}]\label{thm:measureid}
Every stage of the tower is the literal pushforward of the base
measure under the composed map:
$\mu^{(k)} = \bigl(B^{(k)}_{M_0, n}\bigr)_* \mu$.
\end{theorem}

\begin{theorem}[$k$-fold pullback locality;
\lean{kStepBlockObservable\_congr}]\label{thm:congr}
If two fine configurations agree, in positive-edge coordinates, on the
$k$-step pullback support of a stage-$k$ plaquette $P$
(\lean{kStepPullbackSupport} --- the support transport matched
step-by-step to the composed map's own recursion), then the composed
pullback observables agree:
$f(\holo(B^{(k)} A, P)) = f(\holo(B^{(k)} B, P))$.  At $k = 1$ this
recovers the proved one-step locality on the eight-edge support, and
the pullback support object itself agrees with the
\lean{kStepBlockSupport} geometry of \S\ref{sec:support} in every
statement, through a bridge theorem quantified over all predicate
families --- no cast, no \lean{HEq}, no hypothesis
(\lean{kStepPullbackSupport\_succ\_fine},
\lean{kStepPullbackSupport\_terminal\_bridge}); the separation bound
\eqref{eq:termsep} transfers verbatim onto the operational supports
(\lean{terminalPullbackSupport\_separation\_walk},
\lean{terminalPullbackSupport\_separation\_dist}).
\end{theorem}

\begin{theorem}[three-term integration identity;
\lean{truncatedCorrelator\_towerMeasure},
\lean{terminalCorrelator\_eq\_integral\_pullback}]\label{thm:integr}
For measurable $f$, the truncated correlator of the stage-$k$ tower
measure equals the truncated-correlator combination of
\emph{base}-measure integrals of the composed-pullback observables ---
the product integral and both single-observable integrals transported
through the pushforward:
\[
T_{\mu^{(k)}}(p, q)
 = \int f_p^{(k)} f_q^{(k)}\, d\mu
 - \int f_p^{(k)}\, d\mu \int f_q^{(k)}\, d\mu,
\qquad f_p^{(k)}(A) := f\bigl(\holo(B^{(k)} A, p)\bigr),
\]
and, instantiated at the gate's exact objects, the terminal infrared
correlator $c_{\kT(n)}(u)$ equals the three-term base-measure integral
expression at the exact canonical pair.
\end{theorem}

\begin{theorem}[operational support certification;
\lean{terminal\_support\_certified}]\label{thm:capstone}
Fix $d \ge 2$, $N_c \ge 1$, a measurable $f$, $\beta \in \mathbb{R}$,
a base $M_0$, a depth $n \ge 1$, and $u$ in the window
\eqref{eq:window}.  Let $\mu = \muW_{2^n M_0, \beta}$, let
$S_0, S_u$ be the $n$-step pullback supports of the terminal canonical
pair, and let $f^{(n)}_0, f^{(n)}_u$ be the composed-pullback
observables.  Then, jointly:
\begin{enumerate}
\item[(1)] the terminal infrared object consumed by
Definition~\ref{def:gate} equals the base-Wilson-measure truncated
correlator of $f^{(n)}_0$ and $f^{(n)}_u$ (all three integrals
transported);
\item[(2, 3)] each of $f^{(n)}_0, f^{(n)}_u$ depends only on the base
configuration's positive-edge values on $S_0$, resp.\ $S_u$;
\item[(4)] every touch-walk between fine plaquettes whose supports
meet $S_0$ and $S_u$ respectively has length at least
$2^n(2u) - (2^n + 1)$;
\item[(5)] that bound is strictly positive on the whole gate window.
\end{enumerate}
\end{theorem}

Theorem~\ref{thm:capstone} is the sense in which the interface is
\emph{operational}: the infrared clause of the gate is a bound on a
quantity that is \emph{proved} to be a correlation of two explicitly
supported, provably separated functionals of the base Wilson field ---
not on an opaque number attached to a blocked measure.  It certifies
support and metric, not decay: no decay bound and no gate witness is
claimed.

\section{The conditional theorem}\label{sec:main}

\begin{theorem}[conditional physical-distance decay]
\label{thm:main}
Fix $d \ge 2$, $N_c \ge 1$, $f : \SU(N_c) \to \mathbb{R}$, and
$\beta \in \mathbb{R}$.  \textbf{If}
\lean{PhysicalTerminalScaleWilsonGate} holds for these data
(Definition~\ref{def:gate}), \textbf{then} there exist
\[
C \;=\; C_1 + C_2 \sum_{k=0}^{\infty} g_k^{\kappa_0} \;<\; \infty,
\qquad
m \;=\; \min(\varepsilon, c_0) \;>\; 0,
\]
--- \emph{one pair, quantified before every base and depth} --- such
that for every base $M_0 \ge 4$, every depth $n \ge 1$, and every $u$
with $1 \le u$ and $4u \le M_0$,
\[
\Bigl|\, \Cov_{\muW_{2^n M_0,\, \beta}}\bigl(2^n u\bigr) \Bigr|
 \;\le\; C\, e^{-m \cdot 2^n u} .
\]
The proof chain is visible in the formal artifact: the physical
correlator is definitionally the scale-zero tower object
(Proposition~\ref{prop:telescope}); the telescoping ends at the same
fixed index $\kT(n)$ the infrared clause consumes
(Remark~\ref{rem:repair4}); the ultraviolet sum is dominated by
$C_2 e^{-c_0 2^n u} \sum_k g_k^{\kappa_0}$ using the flow summability
(Definition~\ref{def:flow}); and no clause at $k = 0$ is consumed
anywhere --- the bound on the physical object is derived, not
hypothesized.  By Theorem~\ref{thm:capstone}, the infrared quantity
the hypothesis constrains is itself certified as a correlation of two
explicitly supported observables of the base Wilson field at provably
positive, depth-increasing support separation \eqref{eq:termsep}.
\end{theorem}

\begin{remark}[exactly what is claimed]\label{rem:claims}
The Lean artifact of Theorem~\ref{thm:main} is
\lean{wilson\_physical\_mass\_gap\_of\_physicalTerminalGate}.
Theorem~\ref{thm:main} is a conditional statement about a family of
finite-torus correlators.  It is \emph{not} asserted that the gate has
a witness; it is \emph{not} asserted that the Wilson measure satisfies
the (IR)/(UV) clauses at any coupling; the theorem controls the
separations $2^n u$ within the stated windows and no others; and no
thermodynamic or continuum limit is constructed.  The value of the
theorem is architectural: it reduces a physical-distance,
volume-family-uniform decay conclusion to a typed, non-circular,
relabeling-proof, operationally support-certified pair of bound
families --- and Section~\ref{sec:counter} documents why each of those
adjectives costs something.
\end{remark}

\begin{remark}[non-vacuity of the transport machinery]
\label{rem:nonvacuity}
The consumer machinery itself is not vacuous, and this is audited at
the physical measure: in the proved strong-coupling window of the
repository's cluster-expansion lane (the clustering theorem
\lean{sun\_lattice\_exponential\_clustering}; cf.\
\cite{ErikssonAL,OS78}), the \emph{abstract} bridge pipeline of
\S\ref{sec:cm1} fires end to end for the actual $\SU(N_c)$ Wilson
distance-indexed correlator, through the trivial decomposition, and
yields its exponential decay
(\lean{wilson\_correlator\_mass\_gap\_strong\_coupling}).  As the
module states on its face, that witness carries \emph{no} ultraviolet
content --- the strong-coupling decay was already available directly;
it demonstrates only that the transport theorems are inhabited and
fire at the physical measure.  No analogous witness exists for the
terminal gate of Definition~\ref{def:gate}, whose strong-coupling
discharge is blocked precisely by the volume-uniformity and
nonvanishing clauses --- the per-size strong-coupling constants do not
assemble into one pack, and no claim that they do is made.
\end{remark}

\section{The open frontier}\label{sec:frontier}

Everything in this section is open, named, and carried in the
development as definitions without witnesses.  Table~\ref{tab:open}
lists the named open propositions.

\begin{enumerate}
\item \textbf{The gate witness.}  No witness of
\lean{PhysicalTerminalScaleWilsonGate} (or of any of the retained
weaker gates) is provided.  The proposition could in principle be
unsatisfiable; nothing in this paper argues either way.

\item \textbf{The physical Wilson bounds.}  Discharging the (IR) and
(UV) clauses for the actual Wilson measure is the open analytic
mathematics: a single-scale activity/decay estimate for the blocked
gauge theory of Balaban type
\cite{Bal87,Bal88a,Bal88b,Bal89,Dim13a,Dim13b,Dim14}, together with an
infrared bound for the effective theory at the terminal scale.  The
repository's polymer-expansion substrate \cite{ErikssonAL,ErikssonRT}
is theorem-fed toward such estimates, but no part of them is proved
here, and this paper claims no progress on them.

\item \textbf{Odd separations.}  The gate's conclusion family
$\{2^n u\}$ misses every odd separation at every depth $n \ge 1$ ---
an exact theorem (\lean{odd\_separation\_needs\_base\_scale}), not a
gap in exposition.  The named interpolation surrogate
(\lean{CanonicalDoublingDomination}: doubling-window ratio
boundedness) would extend a dyadic bound to all separations at or
above the coarsest covered spacing with explicit adjusted constants
(\lean{allSeparations\_on\_tower\_of\_dyadic\_and\_domination},
proved); but the clause itself is unproved for the Wilson data, is
satisfiable trivially on data with vanishing correlators
(\lean{doubling\_domination\_of\_const}), and admits a formal
obstruction: a zero followed by a nonzero value inside one doubling
window excludes every finite domination constant
(\lean{no\_finite\_domination\_of\_zero\_then\_nonzero},
\lean{canonicalDoublingDomination\_obstruction}).

\item \textbf{Thermodynamic and continuum limits.}  All objects are
finite-torus; the conclusion of Theorem~\ref{thm:main} is uniform over
an unbounded volume family but no infinite-volume correlator is
defined, no limit is shown to exist, and \emph{a fortiori} nothing is
said about a continuum theory.

\item \textbf{Intermediate-remainder locality.}  The operational
certification of \S\ref{sec:oper} covers the terminal stage
$k = \kT(n) = n$.  The supports of the intermediate remainders
$R_k(u)$, $k < n$ --- differences of correlators at consecutive
stages --- are not certified; the composed-pullback locality theorem
for those differences is the natural next lemma.

\item \textbf{Measurability packaging.}  The gate of
Definition~\ref{def:gate} does not carry $\mathrm{Measurable}\, f$ in
its type; Theorem~\ref{thm:capstone} assumes it.  Packaging the
measurability of the observable into the gate type, so that the
certification applies automatically to every admitted instance, is a
small named improvement.
\end{enumerate}

\begin{table}[ht]
\centering
\footnotesize
\begin{tabular}{@{}lll@{}}
\toprule
Named open proposition & Lean name & File \\
\midrule
metric lower transport & \lean{PlaqEmbedDistLowerBound} & \lean{ConcreteGaugeRGFidelity.lean} \\
dyadic $\Rightarrow$ literal uniformity & \lean{DyadicImpliesLiteralQuestion} & \lean{ConcreteGaugeRGFidelity.lean} \\
doubling interpolation clause & \lean{CanonicalDoublingDomination} & \lean{ConcreteGaugeRGLiteralGate.lean} \\
all-separations gate & \lean{AllSeparationsCanonicalWilsonGate} & \lean{ConcreteGaugeRGLiteralGate.lean} \\
terminal gate (renormalized) & \lean{TerminalScaleWilsonGate} & \lean{ConcreteGaugeRGTerminalGate.lean} \\
terminal gate (physical) & \lean{PhysicalTerminalScaleWilsonGate} & \lean{ConcreteGaugeRGPhysicalGate.lean} \\
\bottomrule
\end{tabular}
\caption{Named open propositions: exact definitions, no witnesses, no
satisfiability claims.  The gates are hypotheses of conditional
theorems only.}\label{tab:open}
\end{table}

\medskip
\noindent\textbf{Standing honesty.}  No reduction from anything in
this development to the Clay Yang--Mills problem is claimed, and
importance is not inherited by thematic proximity: only a proved
reduction transfers it.  Within formalized quantum field theory, the
construction of the \emph{free} field with its axioms \cite{DHMN26}
and the present interacting-RG \emph{interface} are, in our view, the
two complementary kinds of artifact the subject currently supports:
constructions that exist, and specifications --- machine-checked,
adversarially hardened --- of what remains to be constructed.

\section{Reproducibility}\label{sec:repro}

Toolchain: Lean \lean{leanprover/lean4:v4.29.0-rc6}; Mathlib pinned to
commit
\begin{center}\small
\lean{07642720480157414db592fa85b626dafb71355b}
\end{center}
\noindent (lakefile and manifest agree); \lean{tectonic} for this
manuscript.

\medskip
\noindent Recorded build facts: at the mathematics freeze commit
\begin{center}\small
\lean{d75d8952} \emph{(short id; this manuscript is added in a
strictly editorial later commit that changes no Lean source)}
\end{center}
\noindent \lean{lake build YangMillsCore} completed successfully
(\textbf{8409 jobs}, zero \lean{sorry}, zero project axioms); the
axiom oracle ran \textbf{2184} \lean{\#print axioms} invocations
(2181 distinct declarations; 3 historical duplicates), every one
printing exactly
\begin{center}
\lean{[propext, Classical.choice, Quot.sound]}
\end{center}
\noindent with zero \lean{sorryAx} and zero nonstandard axioms.  The
oracle transcripts are committed under
\lean{docs/oracle-transcripts/}; the development history, including
the countermodel record of Section~\ref{sec:counter} and every audit
round, is in \lean{docs/C6-BRIDGE-CHARTER.md} --- per the repository's
audit protocol it is archival material and plays no role in the
mathematics above.

\medskip
\noindent Headline-theorem-to-artifact map.  All modules live under
\lean{YangMills/RG/}; the File column abbreviates
\lean{CorrelatorBridge.lean} (Bridge),
\lean{CorrelatorBridgeProvenance.lean} (Provenance),
\lean{ConcreteGaugeRGStep.lean} (Step),
\lean{ConcreteGaugeRGFidelity.lean} (Fidelity),
\lean{ConcreteGaugeRGLiteralGate.lean} (Literal),
\lean{ConcreteGaugeRGTerminalGate.lean} (Terminal),
\lean{ConcreteGaugeRGPhysicalGate.lean} (Physical),
\lean{ConcreteGaugeRGSupport.lean} (Support).  Status \emph{exact} =
in-core Lean theorem or definition with direct oracle witness.

\begin{center}
\scriptsize
\begin{tabular}{@{}llll@{}}
\toprule
Result & Lean name & File & Status \\
\midrule
abstract bridge + transport & \lean{CorrelatorBridge}, \lean{physical\_mass\_gap\_of\_bridge} & Bridge & exact \\
strong-coupling anchor & \lean{wilson\_correlator\_mass\_gap\_strong\_coupling} & Bridge & exact \\
produced remainders (\S\ref{sec:cm1}) & \lean{RGFlowFamily}, \lean{provenant\_decomposition} & Provenance & exact \\
Wilson measure + anchor & \lean{wilsonGibbsMeasure}, \lean{rayCorrelatorOfMeasure\_wilsonGibbs} & Provenance & exact \\
provenance of nonzero UV & \lean{step\_moves\_of\_provenantRsc\_ne\_zero} & Provenance & exact \\
decimation, Prop.~\ref{prop:blockmap} & \lean{blockMap}, \lean{blockMap\_local}, \lean{blockMap\_gaugeAct} & Step & exact \\
effective measure & \lean{effectiveMeasure}, \lean{effectiveMeasure\_isProbability} & Step & exact \\
observable transport & \lean{truncatedCorrelator\_effectiveMeasure} & Step & exact \\
tower + clamp theorems & \lean{towerMeasure}, \lean{lt\_of\_concreteRsc\_ne\_zero} & Step & exact \\
volume-uniform gate + consumer & \lean{NontrivialConcreteRGWilsonBridge} + consumer & Step & exact \\
factor-3 transport, Prop.~\ref{prop:onestep}(i) & \lean{plaqEmbed\_dist\_le\_three\_mul} & Fidelity & exact \\
canonical metric, \eqref{eq:candist} + (ii) & \lean{canonicalPlaquette\_dist\_eq}, \lean{\dots\_dist\_doubles} & Fidelity & exact \\
Wilson probability threading & \lean{wilsonGibbsMeasure\_isProbability} (+ tower) & Fidelity & exact \\
telescoping, Prop.~\ref{prop:telescope} & \lean{scaledCanonical\_decomposition} & Fidelity & exact \\
per-base fidelity gate + consumer & \lean{FidelityConcreteRGWilsonGate} + consumer & Fidelity & exact \\
literal gate + all-bases consumer & \lean{LiteralFidelityConcreteRGWilsonGate} + consumer & Literal & exact \\
odd/small-$t$ diagnosis (\S\ref{sec:cm4}) & \lean{odd\_dyadic\_factorization\_trivial} etc. & Literal & exact \\
depth-zero shell implication & \lean{allSeparations\_of\_literalFidelityGate} & Literal & exact \\
one-step support + separation & \lean{blockPlaquetteSupport}, \lean{blockSupport\_\dots\_dist} & Literal & exact \\
conditional de-lacunarization & \lean{allSeparations\_on\_tower\_of\_dyadic\_and\_domination} & Literal & exact \\
terminal index, Def.~\ref{def:kterm} & \lean{kTerm}, \lean{terminalGate\_IR\_index\_eq}, \lean{rsc\_is\_difference} & Terminal & exact \\
domination obstruction & \lean{no\_finite\_domination\_of\_zero\_then\_nonzero} & Terminal & exact \\
physical gate, Def.~\ref{def:gate} & \lean{PhysicalTerminalScaleWilsonGate} & Physical & exact \\
Theorem~\ref{thm:main} & \lean{wilson\_physical\_mass\_gap\_of\_physicalTerminalGate} & Physical & exact \\
gate comparison & \lean{terminalScaleGate\_of\_physicalTerminalGate} & Physical & exact \\
Theorem~\ref{thm:radius} & \lean{kStepBlockSupport\_source\_in\_ball} & Physical & exact \\
Theorem~\ref{thm:ksep} + positivity & \lean{kStepSupport\_\dots\_dist}, \lean{kStepSeparation\_pos} & Physical & exact \\
terminal tie, \eqref{eq:termsep} & \lean{terminalIR\_is\_canonical\_pair\_correlator} etc. & Physical & exact \\
composed map + Thm.~\ref{thm:measureid} & \lean{iteratedBlockMap}, \lean{towerMeasure\_eq\_map\_iterated} & Support & exact \\
Theorem~\ref{thm:congr} & \lean{kStepBlockObservable\_congr} (+ bridge) & Support & exact \\
Theorem~\ref{thm:integr} & \lean{truncatedCorrelator\_towerMeasure} (+ terminal) & Support & exact \\
Theorem~\ref{thm:capstone} & \lean{terminal\_support\_certified} & Support & exact \\
\bottomrule
\end{tabular}
\end{center}

\medskip
\noindent From a clean clone at the freeze revision (on Windows, clone
with \lean{core.autocrlf=false} and \lean{core.longpaths=true}):
\begin{verbatim}
lake exe cache get
lake build YangMillsCore
lake env lean oracle_check.lean
tectonic -X compile papers/c6-rg-interface/c6_rg_interface.tex
\end{verbatim}

\begin{thebibliography}{99}

\bibitem{Bal87} T.~Ba{\l}aban,
\emph{Renormalization group approach to lattice gauge field theories.
I.~Generation of effective actions in a small field approximation and
a coupling constant renormalization in four dimensions},
Comm. Math. Phys. \textbf{109} (1987), 249--301.
doi:10.1007/BF01215223.

\bibitem{Bal88a} T.~Ba{\l}aban,
\emph{Renormalization group approach to lattice gauge field theories.
II.~Cluster expansions},
Comm. Math. Phys. \textbf{116} (1988), 1--22.
doi:10.1007/BF01239022.

\bibitem{Bal88b} T.~Ba{\l}aban,
\emph{Convergent renormalization expansions for lattice gauge
theories}, Comm. Math. Phys. \textbf{119} (1988), 243--285.
doi:10.1007/BF01217741.

\bibitem{Bal89} T.~Ba{\l}aban,
\emph{Large field renormalization. II.~Localization, exponentiation,
and bounds for the R operation},
Comm. Math. Phys. \textbf{122} (1989), 355--392.
doi:10.1007/BF01238433.

\bibitem{Dim13a} J.~Dimock,
\emph{The renormalization group according to Balaban.
I.~Small fields},
Rev. Math. Phys. \textbf{25} (2013), no.~7, 1330010.
doi:10.1142/S0129055X13300100.  arXiv:1108.1335.

\bibitem{Dim13b} J.~Dimock,
\emph{The renormalization group according to Balaban.
II.~Large fields},
J. Math. Phys. \textbf{54} (2013), 092301.
doi:10.1063/1.4821275.  arXiv:1212.5562.

\bibitem{Dim14} J.~Dimock,
\emph{The renormalization group according to Balaban.
III.~Convergence},
Ann. Henri Poincar\'e \textbf{15} (2014), 2133--2175.
doi:10.1007/s00023-013-0303-3.

\bibitem{OS78} K.~Osterwalder and E.~Seiler,
\emph{Gauge field theories on a lattice},
Ann. Physics \textbf{110} (1978), 440--471.
doi:10.1016/0003-4916(78)90039-8.

\bibitem{Cha19} S.~Chatterjee,
\emph{Yang--Mills for probabilists}, in: Probability and Analysis in
Interacting Physical Systems, Springer Proc. Math. Stat. \textbf{283},
Springer, Cham, 2019, pp.~1--16.
doi:10.1007/978-3-030-15338-0\_1.

\bibitem{Cha26} S.~Chatterjee,
\emph{A scaling limit of $\SU(2)$ lattice Yang--Mills--Higgs theory},
Probab. Math. Phys. \textbf{7} (2026), 339--381.
arXiv:2401.10507.

\bibitem{DHMN26} M.~R.~Douglas, S.~Hoback, A.~Mei, and R.~Nissim,
\emph{Formalization of QFT}, arXiv:2603.15770, 2026.

\bibitem{mathlib20} The mathlib Community,
\emph{The Lean mathematical library}, in: Proc. CPP 2020, ACM, 2020,
pp.~367--381. doi:10.1145/3372885.3373824.

\bibitem{ErikssonAL} L.~Eriksson,
\emph{A machine-checked volume-uniform Wilson-loop area law via a
formalized cluster expansion}, same repository;
ai.viXra.org:2607.0005, July 2026.

\bibitem{ErikssonRT} L.~Eriksson,
\emph{Machine-checked rooted-tree majorants for polymer expansions
with holes}, same repository; ai.viXra.org:2607.0025, July 2026.

\bibitem{Eriksson2D} L.~Eriksson,
\emph{Exact two-dimensional $\SU(2)$ Yang--Mills in Lean}, same
repository; ai.viXra.org:2607.0035, July 2026.

\bibitem{ErikssonC2} L.~Eriksson,
\emph{One Amos bound, three consumer sites: machine-checked
Bessel-ratio calculus for lattice gauge expansions}, companion paper,
same repository, release \lean{c2-v1.2.1};
ai.viXra.org:2607.0033, July 2026.

\bibitem{ErikssonC3} L.~Eriksson,
\emph{A machine-checked proof of the Amos bound for modified Bessel
function ratios}, companion paper, same repository, release
\lean{c3-v1.0.1}; ai.viXra.org:2607.0032, July 2026.

\bibitem{ErikssonC4} L.~Eriksson,
\emph{A machine-checked proof of an Amos-type bound for modified
Bessel ratios at real order}, companion paper, same repository,
release \lean{c4-v1.0.1}; ai.viXra.org:2607.0030, July 2026.

\bibitem{ErikssonC5} L.~Eriksson,
\emph{The diagonal Amos-type family at real order: a machine-checked
quantitative crossing classification}, companion paper, same
repository, release \lean{c5-v1.0.1};
ai.viXra.org:2607.0037, July 2026.

\end{thebibliography}

\end{document}
